% =========================================================
% 6.1 ARQUITECTURA DEL SISTEMA
% El \section{} de esta sección está en el chapter.tex del capítulo;
% sus referencias van en seccion6_1.bib (misma carpeta).
% =========================================================
Para documentar la arquitectura del sistema se utiliza el modelo C4, propuesto por Simon Brown entre 2006 y 2011 \cite{simon_brown_history}. Este modelo describe la arquitectura de software mediante diagramas jerárquicos que van de lo general a lo particular, de modo que cada nivel constituye un acercamiento del anterior. Define cuatro niveles de abstracción: contexto, contenedores, componentes y código\cite{brown_software_architecture}. 

En este trabajo se emplean los dos primeros, dado que estos diagramas buscan el objetivo de comunicar el alcance del sistema, sus actores y las decisiones tecnológicas principales sin descender a detalles de implementación. El diagrama de contexto sitúa el sistema frente a sus usuarios y a los sistemas externos con los que interactúa, mientras que el de contenedores muestra las unidades desplegables que lo componen y la forma en que se comunican. Se omiten los niveles de componentes y código porque la organización interna de cada contenedor responde a decisiones de implementación que cambian con frecuencia, por lo que su representación duplicaría la estructura del código fuente y se desactualizaría rápidamente. Además, el propio Brown señala que los diagramas de contexto y contenedores son suficientes para la mayoría de los equipos \cite{brown_software_architecture}.


\subsection{Diagrama de contexto}
El diagrama de contexto presenta el sistema como una sola caja, junto con los usuarios y los sistemas externos con los que interactúa \cite{brown_software_architecture}. La figura \ref{fig:diagrama_contexto} muestra el diagrama correspondiente al presente proyecto.

% Dibujado con TikZ (no es imagen) para que el texto tenga el tamaño del documento
\begin{sidewaysfigure}
    \centering
    % =========================================================
% Diagrama de contexto (C4, nivel 1)        6.1 Arquitectura del sistema
% ---------------------------------------------------------
% Se incluye con \input dentro del entorno figure de seccion6_1.tex.
% Estilos y colores: setup/settings.tex, sección G (Diagramas C4).
% =========================================================
\begin{tikzpicture}

  % ---------------- Elementos ----------------
  % Persona: caja + cabeza encima
  \node[c4persona, text width=7cm] (usuario)
    {\elementoCuatro{Usuario del sistema}{Persona}};
  \filldraw[c4cabeza] ($(usuario.north)+(0,0.6)$) circle[radius=0.8cm];

  % Sistema
  \node[c4sistema, text width=6cm, below=3cm of usuario] (sistema)
    {\elementoCuatro{Sistema de rutas seguras}{Sistema de software}};

  % Sistemas externos: alineados por su borde superior, 3.8 cm debajo del sistema
  % =================================================
  \node[c4externo, text width=3cm, minimum height=2.7cm, anchor=north] (notificaciones)
    at ($(sistema.south)+(0,-3)$)
    {\elementoCuatro{Servicio de notificaciones push}{Sistema externo}};
  % =================================================
  \node[c4externo, text width=3cm, minimum height=2.7cm, anchor=north] (geocodi)
    at ($(notificaciones.north)+(-4,0)$)
    {\elementoCuatro{Servicio de geocodificación}{Sistema externo}};
  % =================================================
    \node[c4externo, text width=3cm, minimum height=2.7cm, anchor=north] (mapas)
    at ($(notificaciones.north)+(-8,0)$)
    {\elementoCuatro{Servicio de mapas}{Sistema externo}};
 % =================================================    
  \node[c4externo, text width=3cm, minimum height=2.7cm, anchor=north] (google)
    at ($(notificaciones.north)+(4,0)$)
    {\elementoCuatro{Identidad de Google}{Sistema externo}};
 % =================================================
  \node[c4externo, text width=3cm, minimum height=2.7cm, anchor=north] (email)
    at ($(notificaciones.north)+(8,0)$)
    {\elementoCuatro{Servicio de email}{Sistema externo}};

  % ---------------- Relaciones ----------------
  \draw[c4relacion] (usuario) -- (sistema)
    node[c4etiqueta, midway, text width=7.4cm]
    {Gestiona su perfil y favoritos, consulta mapas y rutas seguras,
     genera/valida reportes y comparte o monitorea ubicaciones usando};

  % --- Fila inferior: tres flechas verticales rectas (x = -3, 0, +3) ---
  % Cada x cae dentro del ancho de su caja destino (cajas de 3 cm).
  \draw[c4relacion] ($(sistema.south)+(-3,0)$) -- ($(geocodi.north)+(1,0)$)
    node[c4etiqueta, midway, text width=2.5cm]
    {Busca lugares por nombre o dirección usando};

  \draw[c4relacion] (sistema.south) -- (notificaciones.north)
    node[c4etiqueta, midway, text width=2.5cm]
    {Envía alertas al usuario usando};

  \draw[c4relacion] ($(sistema.south)+(3,0)$) -- ($(google.north)+(-1,0)$)
    node[c4etiqueta, midway, text width=2.5cm]
    {Autentica a los usuarios usando};

  % --- Cajas laterales (mapas y email): líneas rectas desde las esquinas inferiores del sistema ---
  \draw[c4relacion] (sistema.south west) -- (mapas.north)
    node[c4etiqueta, midway, text width=2.3cm]
    {Obtiene los mapas base de};

  \draw[c4relacion] (sistema.south east) -- (email.north)
    node[c4etiqueta, midway, text width=2.3cm]
    {Envía correos electrónicos usando};

\end{tikzpicture}
    \caption{Diagrama de contexto del sistema}
    \label{fig:diagrama_contexto}
\end{sidewaysfigure}

\subsection{Diagrama de contenedores}
El diagrama de contenedores descompone el sistema en las aplicaciones y los almacenes de datos que lo conforman, e indica la tecnología utilizada y la comunicación entre ellos \cite{brown_software_architecture}. La figura \ref{fig:diagrama_contenedor} muestra el diagrama correspondiente al presente proyecto.
% Dibujado con TikZ (no es imagen) para que el texto tenga el tamaño del documento
\begin{sidewaysfigure}
    \centering
    % =========================================================
% Diagrama de contenedores (C4, nivel 2)
% =========================================================
\begin{tikzpicture}

  % ---------------- Contenedores ----------------
  \node[c4contenedor, text width=4.5cm] (auth) at (1.5, -5)
    {\elementoCuatro{Autenticación}{Contenedor: Supabase}};

  \node[c4contenedor, text width=4.5cm] (app) at (4.5, 0)
    {\elementoCuatro{Aplicación móvil}{Contenedor: Flutter}};

  \node[c4basedatos, text width=4.5cm] (db) at (4.5, -10)
    {\elementoCuatro{Base de datos}{Contenedor: PostgreSQL}};

  \node[c4contenedor, text width=4.5cm] (api) at (7.5, -5)
    {\elementoCuatro{API REST}{Contenedor: Supabase}};

  % ---------------- Límite del sistema ----------------
  \node[c4limiteEtiqueta, anchor=north west] (etiquetaLimite)
    at (-2, 1.5)
    {\textbf{Sistema de rutas seguras}\\{[Sistema de software]}};

  \begin{scope}[on background layer]
    \node[c4limite, fit=(etiquetaLimite)(api)(app)(db)(auth)] (limite) {};
  \end{scope}

  % ---------------- Persona ----------------
  \node[c4persona, text width=6cm] (usuario) at (4.5, 3.5)
    {\elementoCuatro{Usuario del sistema}{Persona}};
  \filldraw[c4cabeza] ($(usuario.north)+(0,0.6)$) circle[radius=0.8cm];

  % ---------------- Sistemas externos ----------------
  \node[c4externo, text width=2.8cm, minimum height=2.7cm, anchor=north] (notificaciones)
    at ($(api.north)+(6,-4)$)
    {\elementoCuatro{Servicio de notificaciones push}{Sistema externo}};

  \node[c4externo, text width=2.8cm, minimum height=2.7cm, anchor=north] (email)
    at ($(auth.north)+(-6,-4)$)
    {\elementoCuatro{Servicio de email}{Sistema externo}};

  \node[c4externo, text width=2.8cm, minimum height=2.7cm, anchor=north] (google)
    at ($(auth.north)+(-6,0)$)
    {\elementoCuatro{Identidad de Google}{Sistema externo}};

  \node[c4externo, text width=2.8cm, minimum height=2.7cm, anchor=north] (geocodi)
    at ($(api.north)+(6,1.5)$)
    {\elementoCuatro{Servicio de geocodificación}{Sistema externo}};

  \node[c4externo, text width=2.8cm, minimum height=2.7cm, anchor=north] (mapas)
    at ($(api.north)+(6,5)$)
    {\elementoCuatro{Servicio de mapas}{Sistema externo}};

  % ---------------- Relaciones Internas ----------------
  \draw[c4relacion] (app.south) -- (auth.north);
  \draw[c4relacion] (app.south) -- (api.north);
  \draw[c4relacion] (auth.south) -- (db.north);
  \draw[c4relacion] (api.south) -- (db.north);

  % ---------------- Relaciones Externas y Usuario ----------------
  \draw[c4relacion] (usuario) |- (app)
    node[c4etiqueta, midway, right] {Interactúa con};

  \draw[c4relacion] (email.west) |- ++(-0.5,0) |- (usuario.west)
    node[c4etiqueta, pos=0.25, right, text width=2.5cm] {Envía email de verificación};

  \draw[c4relacion] (auth.south) |- (email.east)
    node[c4etiqueta, near end, above] {Envía verificación};

  \draw[c4relacion] (auth.north) |- (google.east)
    node[c4etiqueta, near end, below] {Verifica/agrega token};

  \draw[c4relacion] (app.west) -- (google.north)
    node[c4etiqueta, pos=0.25, above] {Obtiene token de Google};

  \draw[c4relacion] (app.east) |- (mapas.west)
    node[c4etiqueta, midway, above] {Obtiene mapas de};

  \draw[c4relacion] (app.east) |- (geocodi.west)
    node[c4etiqueta, near end, above] {Busca lugares usando};

  \draw[c4relacion] (api.south) |- (notificaciones.west)
    node[c4etiqueta, near end, above] {Envía alertas vía};

\end{tikzpicture}
    \caption{Diagrama de contenedor del sistema}
    \label{fig:diagrama_contenedor}
\end{sidewaysfigure}

\clearpage
\restoregeometry